% ============================================
% LUC-08
% Gerado de molambudos.md — build_miolo.py
% ============================================
\section*{\textbf{LUC-08} --- A Última Sessão de Lúcia}
\markboth{LUC-08 --- A Última Sessão de Lúcia}{}

\noindent \textit{A psiquiatra vira paciente}

\noindent \textbf{Maio de 2026 --- Consultório de Lúcia, Belo Horizonte}

\vspace{0.5em}\noindent\rule{\textwidth}{0.5pt}\vspace{0.5em}

\noindent Ela guardou tudo na gaveta.
\noindent Pelo menos foi isso que declarou depois.

\noindent Não como profissional --- como paciente.

\noindent A Dra. Regina Alves era amiga de Lúcia desde a faculdade. Atendia num consultório elegante em BH, com poltronas de couro e uma vista para a serra. Lúcia sentou na poltrona do paciente pela primeira vez em 15 anos.
\noindent \textquotedblleft{\textit{O que está acontecendo, Lúcia?\textquotedblright{}}}
\noindent \textquotedblleft{\textit{Eu estou com medo de estar ficando louca.\textquotedblright{}}}

\noindent \textquotedblleft{\textit{O que você acha que está acontecendo?\textquotedblright{}}}
\noindent \textquotedblleft{\textit{Acho que alguém montou um arquivo para produzir uma reação específica em quem lesse. Ou acho que o arquivo encontrou uma forma de continuar sendo lido. As duas hipóteses me assustam do mesmo jeito.\textquotedblright{}}}

\noindent E alguma coisa tinha acabado de encostar nela.

\noindent \textit{\textquotedblleft{}Por quê?\textquotedblright{}}

\noindent \textit{\textquotedblleft{}Estou tendo sintomas psicóticos.\textquotedblright{}}

\noindent \textit{\textquotedblleft{}Quais?\textquotedblright{}}

\noindent Lúcia listou: o cheiro, os sonhos, o olho amarelo, a sensação de não estar sozinha, a letra da mão mudando.

\noindent Regina ouviu em silêncio. Fez anotações.

\noindent Depois de 20 minutos, ela perguntou:

\noindent \textit{\textquotedblleft{}Tem algo mais?\textquotedblright{}}

\noindent \textit{\textquotedblleft{}...Tem um diário.\textquotedblright{}}

\noindent Ela contou sobre o caso do paciente 1.260. Sobre a criatura. Sobre o ciclo.

\noindent Quando terminou, Regina fechou o bloco.

\noindent \textit{\textquotedblleft{}Lúcia, você está com um quadro de estresse agudo. Possivelmente um transtorno dissociativo desencadeado pelo contato com material traumático.\textquotedblright{}}

\noindent \textit{\textquotedblleft{}Você acha que é sugestão?\textquotedblright{}}

\noindent \textit{\textquotedblleft{}Acho que você leu um caso perturbador, seu cérebro processou o trauma, e agora está produzindo sintomas compatíveis com o que você leu. É um mecanismo de identificação.\textquotedblright{}}

\noindent \textit{\textquotedblleft{}E o olho?\textquotedblright{}}

\noindent \textit{\textquotedblleft{}O que tem o olho?\textquotedblright{}}

\noindent Lúcia tirou os óculos escuros.

\noindent O olho direito estava amarelo.

\noindent Regina inclinou a cabeça.

\noindent \textit{\textquotedblleft{}Isso... é novo?\textquotedblright{}}

\noindent \textit{\textquotedblleft{}Apareceu depois que li o diário.\textquotedblright{}}

\noindent \textit{\textquotedblleft{}Você já foi ao oftalmologista?\textquotedblright{}}

\noindent \textit{\textquotedblleft{}Fui. Ele disse que não há explicação.\textquotedblright{}}

\noindent Regina ficou em silêncio.

\noindent \textit{\textquotedblleft{}Lúcia, você acredita que uma história pode causar uma mudança fisiológica no olho?\textquotedblright{}}
\noindent \textquotedblleft{\textit{Eu acredito que exposição repetida a material traumático pode produzir sintomas reais.\textquotedblright{}}}
\noindent Lúcia preferiu acreditar que era a lanterna --- mesmo depois de 47 anos.
\noindent \textquotedblleft{\textit{É a única resposta honesta que eu tenho.\textquotedblright{}}}

\noindent \textit{\textquotedblleft{}Não.\textquotedblright{}}

\noindent \textit{\textquotedblleft{}Mas está acontecendo.\textquotedblright{}}

\noindent \textit{\textquotedblleft{}Está.\textquotedblright{}}

\noindent \textit{\textquotedblleft{}O que você acha que é?\textquotedblright{}}

\noindent Lúcia respirou fundo.

\noindent \textit{\textquotedblleft{}Eu acho que o que está no diário é real. Eu acho que a criatura existe. Eu acho que ela passou de Joaquim para Oliveira. E de Oliveira para mim.\textquotedblright{}}

\noindent Regina não respondeu.

\noindent Depois de um longo silêncio, ela disse:

\noindent \textit{\textquotedblleft{}Lúcia, como amiga, eu te peço: procura um tratamento. Internação, se necessário.\textquotedblright{}}

\noindent \textit{\textquotedblleft{}E como psicóloga?\textquotedblright{}}

\noindent \textit{\textquotedblleft{}Como psicóloga...\textquotedblright{}\textit{ Regina hesitou. }\textquotedblleft{}Como psicóloga, eu te pergunto: o que você vai fazer com esse diagnóstico?\textquotedblright{}}

\noindent Lúcia levantou.

\noindent \textit{\textquotedblleft{}Vou escrever. Vou registrar tudo. Para que, quando o ciclo chegar na próxima pessoa, ela saiba o que está acontecendo.\textquotedblright{}}

\noindent \textit{\textquotedblleft{}E se a próxima pessoa for eu?\textquotedblright{}}

\noindent Lúcia parou na porta.

\noindent \textit{\textquotedblleft{}Então você vai entender.\textquotedblright{}}

\noindent Ela saiu.

\noindent Nunca mais falou com Regina.

\noindent Três semanas depois, recebeu um envelope pelo correio. Sem bilhete. Sem remetente.
\noindent Dentro havia uma chave pequena, ainda úmida, e uma folha arrancada de um diário que Regina nunca tinha visto.

\noindent Lúcia tirou da bolsa um envelope lacrado.
\noindent \textquotedblleft{\textit{Vou entregar o original para custódia judicial e manter só transcrições parciais. Ninguém precisa tocar no diário de novo.\textquotedblright{}}}
\noindent \textquotedblleft{\textit{Isso é uma decisão clínica ou superstição?\textquotedblright{}}}
\noindent Na porta, antes de sair, deixou uma cópia do laudo sobre a mesa de Regina.
\noindent \textquotedblleft{\textit{Se eu desaparecer, entregue isto ao Conselho. Não leia o diário. Leia só o laudo.\textquotedblright{}}}

\noindent Não abriu.

\noindent Não leu.

\noindent Até hoje.

\vspace{0.5em}
\noindent\textit{\textbf{↪ Rotas:} \rota{LUC-09} \rotasep \rota{CONT-04}}
\noindent Ela marcou uma sessão com uma psicóloga.